\section*{ЗАКЛЮЧЕНИЕ}
\phantomsection
\addcontentsline{toc}{section}{Заключение}

В рамках настоящей курсовой работы выполнено численное исследование влияния регулярного дискретного микрорельефа на трибологические характеристики радиального подшипника скольжения. По итогам решены следующие задачи.

\begin{enumerate}
    \item Систематизированы теоретические представления о формировании несущего смазочного слоя и обобщены опубликованные результаты по лазерному текстурированию поверхностей трения. Показано, что определяющий вклад в повышение характеристик подшипника вносит микрогидродинамический эффект, формирующийся за счёт локальных клиновых зазоров на кромках микроуглублений.

    \item Реализована численная схема решения уравнения Рейнольдса для подшипника конечной длины: метод конечных разностей с итерационной процедурой Гаусса--Зейделя и верхней релаксацией. Решение получено на сетке $500 \times 500$ узлов с допуском по невязке $\delta < 10^{-5}$.

    \item Получены поля давления и значения интегральных показателей -- подъёмной силы $F$, коэффициента трения $\mu$ и объёмного расхода $Q$ -- как для гладкого подшипника, так и для подшипника с коническими углублениями (тип~6). Расчёт проведён в диапазоне относительных эксцентриситетов $\varepsilon = 0{,}05\text{ - }0{,}80$ для подшипника геометрии~A: $R = 35$~мм, $c = 50$~мкм, $L = 56$~мм.

    \item Выполнено сопоставление гладкого и текстурированного вариантов. При характерном рабочем значении $\varepsilon = 0{,}6$ зафиксированы следующие количественные изменения:
    \begin{itemize}
        \item подъёмная сила выросла на $2{,}90$\% -- с $F = 5234{,}7$~Н до $F = 5386{,}6$~Н;
        \item коэффициент трения уменьшился на $4{,}06$\% -- с $\mu = 0{,}01480$ до $\mu = 0{,}01420$;
        \item объёмный расход смазки увеличился на $0{,}71$\% -- с $Q = 0{,}2356$~л/с до $Q = 0{,}2373$~л/с.
    \end{itemize}
\end{enumerate}

Полученные данные подтверждают, что нанесение конического микрорельефа с параметрами $a = 2{,}35$~мм, $b = 2{,}15$~мм и $h_p = 10$~мкм при относительной глубине $h_p/c = 0{,}2$ оказывает положительное, хотя и умеренное, влияние на рабочие характеристики подшипника. Одновременное увеличение подъёмной силы и снижение коэффициента трения при пренебрежимо малом приросте расхода смазки свидетельствует о принципиальной работоспособности данного типа микрорельефа.

Конический профиль характеризуется монотонно убывающей глубиной от центра углубления к периферии: глубина изменяется линейно по радиусу, а в краевой точке обращается в нуль. Такая геометрия обеспечивает плавное взаимодействие смазочного слоя с микрорельефом и отсутствие резких возмущений потока. Однако линейное изменение глубины формирует менее выраженный микроклин по сравнению с профилями, имеющими резкую границу или участок постоянной глубины, поэтому прирост микрогидродинамического давления оказывается умеренным.

С учётом полученных результатов можно констатировать, что конический микрорельеф является технологически простым и устойчивым к локальным повреждениям решением, однако по эффективности уступает профилям с более резкой клиновой геометрией. Применение конического текстурирования целесообразно в узлах со значительным разбросом эксплуатационных режимов, где приоритетной является устойчивость гидродинамического режима, а не максимальное приращение нагрузочной способности.
